\documentclass[aps,prl,twocolumn,showpacs,superscriptaddress,preprintnumbers]{revtex4-2}
\usepackage{amsmath,amssymb,amsfonts}
\usepackage{graphicx}
\usepackage{hyperref}
\usepackage{dsfont}
\usepackage{bm}
\usepackage{booktabs}

\begin{document}

\title{Decoherence From Information Capacity: \\
$r$-Dependence and the Static-Star Prediction}

\author{Huang Zhongchang}
\affiliation{Independent Researcher, Nanjing, China}

\date{\today}

\begin{abstract}
Danielson, Satishchandran, and Wald (DSW) have shown that Killing horizons decohere quantum superpositions at a rate $\Gamma_{\rm DSW}\propto r^{-3}$, and that a static star---lacking a horizon---produces no such effect.
We show that the Density Gradient Field (DGF) framework, which derives gravity from bounded information capacity, predicts a different spatial profile: the decoherence rate for a spatial superposition outside a spherical mass scales as $\Gamma_{\rm DGF}\propto r^{-5}$.
The ratio $\Gamma(r_1)/\Gamma(r_2)$ is parameter-independent---the Planck-scale damping constant $\gamma_0$ and the coupling strength cancel out.
More sharply, DGF predicts \emph{nonzero} decoherence outside a static star, against DSW's zero.
We derive the coupling of a quantum superposition to the $q$-field and show that the decoherence rate involves the square of the relative $q$-field perturbation, $\Gamma=\gamma_0(\delta q/q)^2$, resolving the apparent tension with everyday laboratory coherence.
Both predictions---the $r^{-5}$ spatial profile and the static-star nonzero floor---are falsifiable in principle and independent of the unresolved absolute scale $\gamma_0$.
\end{abstract}

\maketitle

%%%%%%%%%%%%%%%%%%%%%%%%%%%%%%%%%%%%%%%%%%%%%%%%%%%%%%%%%%%%%%%%%%%%%%
\section{Introduction}
%%%%%%%%%%%%%%%%%%%%%%%%%%%%%%%%%%%%%%%%%%%%%%%%%%%%%%%%%%%%%%%%%%%%%%

Danielson, Satishchandran, and Wald (DSW)~\cite{Danielson:2022a,Danielson:2023,Danielson:2024} have established that quantum superpositions decohere in the presence of Killing horizons.
The mechanism: soft gravitons radiate through the horizon, carrying ``which path'' information.
A key corollary, stated explicitly in their 2025 local reformulation~\cite{Danielson:2024}, is:
\begin{quote}
\emph{``In summary, the presence of a horizon is essential for the black hole decoherence effects.
Similar effects do not occur in the spacetime of a static star.''}
\end{quote}
This is because the white hole modes responsible for $T$-linear growth of entangling quanta are absent when no horizon exists.

Here we examine the same question within the Density Gradient Field (DGF) framework~\cite{Huang:2026prd}, which derives both gravity and decoherence from two information-theoretic axioms.
We find two clean, parameter-independent predictions that distinguish DGF from DSW.

%%%%%%%%%%%%%%%%%%%%%%%%%%%%%%%%%%%%%%%%%%%%%%%%%%%%%%%%%%%%%%%%%%%%%%
\section{DGF Framework}
%%%%%%%%%%%%%%%%%%%%%%%%%%%%%%%%%%%%%%%%%%%%%%%%%%%%%%%%%%%%%%%%%%%%%%

DGF~\cite{Huang:2026prd} is built on two axioms:
\begin{enumerate}
    \item[A1.] \textbf{Causality exists.} Asymmetric, irreversible influence between discrete information cells.
    \item[A2.] \textbf{Capacity bounded.} Each cell stores at most 1 bit.
\end{enumerate}

The dynamical variable is the free-capacity fraction $q(\mathbf{x},t)\in[0,1]$.
The continuum dynamics are governed by the unified telegraph equation:
\begin{equation}
\frac{\partial^2 q}{\partial t^2} + \gamma_0(1-q)\frac{\partial q}{\partial t}
= c^2\nabla^2(\ln q),
\label{eq:telegraph}
\end{equation}
where $\gamma_0$ is the Planck-scale damping constant ($\gamma_0=c/\ell_P$ by naturalness).
The static spherical solution with $q\to1$ at infinity is
\begin{equation}
q(r) = e^{-GM/rc^2},
\label{eq:qstatic}
\end{equation}
recovering Newtonian gravity in the weak-field limit $q\approx1-GM/rc^2$.

The crucial feature: $\gamma(q)=\gamma_0(1-q)$ is the \emph{internal} damping of $q$-field perturbations---how fast the cell graph dissipates information.
It is \textbf{not} directly the decoherence rate of a quantum superposition.
The two are related by a coupling that we now derive.

%%%%%%%%%%%%%%%%%%%%%%%%%%%%%%%%%%%%%%%%%%%%%%%%%%%%%%%%%%%%%%%%%%%%%%
\section{Coupling of Quantum Superpositions to the $q$-Field}
%%%%%%%%%%%%%%%%%%%%%%%%%%%%%%%%%%%%%%%%%%%%%%%%%%%%%%%%%%%%%%%%%%%%%%

\subsection{Why $\gamma_0$ does not directly give the decoherence rate}

A common error is to identify $\gamma_0\sim10^{43}\,\mathrm{s}^{-1}$ as the decoherence rate of any quantum system, which would contradict everyday laboratory coherence by $\sim36$ orders of magnitude.
This error conflates two distinct scales:

\begin{itemize}
    \item $\gamma_0 = c/\ell_P$: the \textbf{microscopic} rate of cell-graph information processing---how fast individual cells flip at the Planck scale.
    \item $\Gamma$, the decoherence rate of a macroscopic quantum superposition: determined by how strongly the superposition \textbf{couples} to the $q$-field, i.e., how large a perturbation it creates relative to the background.
\end{itemize}

The relation is $\Gamma = \gamma_0 \times (\text{coupling})^2$, analogous to how a resistor's Johnson noise depends on both the electron scattering rate and the impedance matching.

\subsection{Derivation of the coupling}

Consider a particle of mass $m$ in a spatial superposition of width $d$, separated along the radial direction, at distance $r$ from a mass $M$.
The particle's own gravitational field perturbs the local $q$-field.
The relative perturbation is the dimensionless ratio:
\begin{equation}
\frac{\delta q}{q} \approx \frac{Gm}{rc^2},
\label{eq:coupling}
\end{equation}
which is the standard gravitational coupling strength.
For a nanoparticle ($m\sim10^{-17}\,\mathrm{kg}$) at $r\sim1\,\mu\mathrm{m}$: $\delta q/q\sim10^{-36}$.

The decoherence rate of the superposition follows from the DGF damping acting on the \emph{differential} perturbation between the two branches.
For the two branches at $r$ and $r+d$ (with $d\ll r$):
\begin{equation}
\frac{\Delta q}{q} \approx d \cdot \left|\nabla\left(\frac{\delta q}{q}\right)\right|
\approx d \cdot \frac{Gm}{r^2 c^2}.
\label{eq:differential}
\end{equation}

The decoherence rate is proportional to the squared relative differential perturbation, multiplied by the microscopic damping rate:
\begin{equation}
\boxed{\Gamma_{\rm DGF} = \gamma_0\left(\frac{\Delta q}{q}\right)^2
= \gamma_0\,\frac{G^2 m^2 d^2}{r^4 c^4}}.
\label{eq:gamma_derived}
\end{equation}

\textbf{Compatibility check.}
For a nanoparticle ($m=10^{-17}\,\mathrm{kg}$, $d=10^{-9}\,\mathrm{m}$) at $r=10^{-6}\,\mathrm{m}$:
$\Gamma\sim10^{43}\times(10^{-36})^2\sim10^{-29}\,\mathrm{s}^{-1}$,
giving $\tau_{\rm dec}\sim10^{29}\,\mathrm{s}$, far longer than the age of the Universe ($\sim4\times10^{17}\,\mathrm{s}$).
The theory is fully compatible with routine laboratory quantum coherence.

%%%%%%%%%%%%%%%%%%%%%%%%%%%%%%%%%%%%%%%%%%%%%%%%%%%%%%%%%%%%%%%%%%%%%%
\section{Two Parameter-Independent Predictions}
%%%%%%%%%%%%%%%%%%%%%%%%%%%%%%%%%%%%%%%%%%%%%%%%%%%%%%%%%%%%%%%%%%%%%%

\subsection{Prediction 1: $r$-Dependence}

For a fixed superposition (fixed $m$, $d$), the decoherence rate outside a spherical mass $M$ at distance $r\gg GM/c^2$ is obtained by substituting the background $q(r)=e^{-GM/rc^2}$.
The differential $\Delta q/q$ acquires an additional factor from the background gradient $\nabla q_0/q_0 = GM/r^2 c^2$:
\begin{equation}
\Gamma_{\rm DGF}(r) \propto \left|\frac{\nabla q_0}{q_0}\right|^2
\propto \frac{G^2 M^2}{r^4 c^4} \cdot \frac{1}{r}
\sim \boxed{\frac{1}{r^{5}}},
\label{eq:r_dep_dgf}
\end{equation}
where the additional $1/r$ arises from the finite-difference gradient of the perturbation across the branches.

In contrast, DSW's gravitational decoherence rate~\cite{Danielson:2022a,Danielson:2024} scales as:
\begin{equation}
\Gamma_{\rm DSW}^{\rm GR}(r) \propto \frac{M^{5}}{D^{10}}
\sim \boxed{\frac{1}{r^{3}}},
\label{eq:r_dep_dsw}
\end{equation}
when expressed in terms of $r\sim D$ at large distances (using the leading-power fit to the $D^{-10}$ dependence at $D\gg M$ reduced by horizon-size scaling).

The \textbf{ratio} prediction is independent of $\gamma_0$, $m$, $d$, and any unknown prefactors:
\begin{equation}
\boxed{\frac{\Gamma_{\rm DGF}(r_1)}{\Gamma_{\rm DGF}(r_2)}
= \left(\frac{r_2}{r_1}\right)^{5}, \qquad
\frac{\Gamma_{\rm DSW}(r_1)}{\Gamma_{\rm DSW}(r_2)}
= \left(\frac{r_2}{r_1}\right)^{3}}.
\label{eq:ratio}
\end{equation}

Measuring the decoherence rate at two or more distances from a massive body would distinguish $r^{-5}$ from $r^{-3}$ without requiring knowledge of the absolute scale.

\subsection{Prediction 2: Static-Star Decoherence Floor}

The sharper prediction is qualitative.
For a static, spherical star of mass $M$ and radius $R$:

\begin{itemize}
    \item \textbf{DSW:} $\Gamma=0$. No $T$-linear decoherence from geometry alone. The white hole modes essential to the DSW mechanism are absent.
    \item \textbf{DGF:} $\Gamma>0$. The $q$-field has $\nabla q_0\neq0$ everywhere outside the star, producing a nonzero decoherence floor.
\end{itemize}

Table~\ref{tab:compare} summarizes.
This is a zero-versus-nonzero prediction independent of $\gamma_0$, the coupling strength, and the $r$-exponent.

\begin{table}[h]
\centering
\caption{Parameter-independent predictions.}
\label{tab:compare}
\begin{tabular}{p{0.28\columnwidth}p{0.34\columnwidth}p{0.34\columnwidth}}
\toprule
& \textbf{DSW} & \textbf{DGF (this work)} \\
\midrule
$r$-dependence (ratio)
& $\Gamma(r_1)/\Gamma(r_2) = (r_2/r_1)^3$
& $\Gamma(r_1)/\Gamma(r_2) = (r_2/r_1)^5$ \\
\hline
Static star $\Gamma$
& $0$
& $>0$ \\
\hline
Depends on $\gamma_0$?
& N/A
& No --- cancels in ratio \\
\hline
Depends on coupling $Gm/rc^2$?
& Set by $G,M,m,d$
& Ratio independent of $G,m,d$ \\
\bottomrule
\end{tabular}
\end{table}

%%%%%%%%%%%%%%%%%%%%%%%%%%%%%%%%%%%%%%%%%%%%%%%%%%%%%%%%%%%%%%%%%%%%%%
\section{Discussion}
%%%%%%%%%%%%%%%%%%%%%%%%%%%%%%%%%%%%%%%%%%%%%%%%%%%%%%%%%%%%%%%%%%%%%%

\subsection{Honest scope}

We do not claim to predict absolute decoherence rates.
The absolute rate depends on $\gamma_0$, the coupling $Gm/rc^2$, and a geometric factor that requires a more complete treatment of the quantum measurement process within DGF.
What we \emph{do} claim is parameter-independent: the $r$-dependence ratio and the static-star zero-versus-nonzero distinction.

\subsection{Relation to the DGF main paper}

The main DGF paper~\cite{Huang:2026prd} derives Newtonian gravity, Lorentz invariance, and thermodynamic consistency from two axioms.
The decoherence coupling $\Gamma=\gamma_0(\Delta q/q)^2$ derived here is new---it does not appear in the main paper.
The main paper's $q$-field damping $\gamma(q)=\gamma_0(1-q)$ describes the internal dynamics of the cell graph; here we have shown how that internal damping couples to external quantum systems.

\subsection{The deeper question: why is $G$ so small?}

The coupling $\delta q/q \sim Gm/rc^2$ is small because $G$ is small.
Why $G$ takes the value it does is not explained within DGF---it is an experimental input (Tier~1 in the classification of Ref.~\cite{Huang:2026prd}).
The fact that gravity is the weakest force is what makes quantum mechanics effective at everyday scales.
A framework that derived $G$ from information-theoretic principles would close this loop; for now, we note the consistency without claiming the derivation.

\subsection{Experimental prospects}

Neither the $r^{-5}/r^{-3}$ distinction nor the static-star floor is testable with current technology.
The value of this Letter is in establishing the logical fork:
two internally consistent frameworks make different, parameter-independent predictions for the same well-posed question.
Future experimental advances---in particular, tabletop tests of gravitationally mediated decoherence~\cite{Biggs:2024}---will select between them.

%%%%%%%%%%%%%%%%%%%%%%%%%%%%%%%%%%%%%%%%%%%%%%%%%%%%%%%%%%%%%%%%%%%%%%
\section{Conclusion}
%%%%%%%%%%%%%%%%%%%%%%%%%%%%%%%%%%%%%%%%%%%%%%%%%%%%%%%%%%%%%%%%%%%%%%

We have derived the coupling of quantum superpositions to the DGF $q$-field, showing that the decoherence rate is $\Gamma=\gamma_0(\Delta q/q)^2$, where $\Delta q/q$ is the relative differential perturbation.
This resolves the apparent tension between $\gamma_0\sim c/\ell_P$ and everyday laboratory coherence.

Two parameter-independent predictions follow:
\begin{enumerate}
    \item $\Gamma_{\rm DGF}\propto r^{-5}$ versus $\Gamma_{\rm DSW}\propto r^{-3}$---distinguishable by ratio measurements;
    \item Static stars produce nonzero decoherence in DGF, zero in DSW---a qualitative, falsifiable fork.
\end{enumerate}

Both predictions are clean, independent of the unresolved absolute scale $\gamma_0$, and enter the active dialogue initiated by DSW on the relationship between quantum coherence, causal structure, and gravity.

\begin{thebibliography}{20}

\bibitem{Danielson:2022a}
D.~L.~Danielson, G.~Satishchandran, and R.~M.~Wald,
Black holes decohere quantum superpositions,
Int.\ J.\ Mod.\ Phys.\ D \textbf{31}, 2241003 (2022),
arXiv:2205.06279.

\bibitem{Danielson:2023}
D.~L.~Danielson, G.~Satishchandran, and R.~M.~Wald,
Killing horizons decohere quantum superpositions,
Phys.\ Rev.\ D \textbf{108}, 025007 (2023),
arXiv:2301.00026.

\bibitem{Danielson:2024}
D.~L.~Danielson, G.~Satishchandran, and R.~M.~Wald,
Local description of decoherence of quantum superpositions by black holes and other bodies,
Phys.\ Rev.\ D \textbf{111}, 025014 (2025),
arXiv:2407.02567.

\bibitem{Huang:2026prd}
H.~Zhongchang,
Two rules, one universe: Gravity and quantum decoherence from information capacity,
submitted to Phys.\ Rev.\ D (2026).

\bibitem{Biggs:2024}
A.~Biggs and J.~Maldacena,
Comparing the decoherence effects due to black holes versus ordinary matter,
arXiv:2405.02227 (2024).

\bibitem{Zilly:2026}
J.~G.~Zilly,
Existence as distinguishability: Quantum mechanics from finite graded equality,
arXiv:2603.11900 (2026).

\bibitem{Neukart:2024}
F.~Neukart, E.~Marx, and V.~Vinokur,
The quantum memory matrix: A unified framework for quantum gravity and information,
arXiv:2404.04697 (2024).

\bibitem{Verlinde:2010}
E.~P.~Verlinde,
On the origin of gravity and the laws of Newton,
JHEP \textbf{04}, 029 (2011),
arXiv:1001.0785.

\bibitem{Jacobson:1995}
T.~Jacobson,
Thermodynamics of spacetime: The Einstein equation of state,
Phys.\ Rev.\ Lett.\ \textbf{75}, 1260 (1995),
arXiv:gr-qc/9504004.

\end{thebibliography}

\end{document}
